\documentclass[11pt]{article}
\usepackage[margin=1in]{geometry}
\usepackage{times}
\usepackage{hyperref}
\hypersetup{colorlinks=true, linkcolor=black, urlcolor=blue, citecolor=black}
\title{Gravity Has No Off Switch: A Survey of Patent-Office Antigravity Claims}
\author{Flyxion \\ \small Independent Researcher}
\date{}
\begin{document}
\maketitle

\begin{abstract}
Patent offices across multiple jurisdictions have, over the past century, granted patents for devices described as gravity shielding, gravity generating, or antigravity propulsion apparatus, despite the patent examination process not requiring proof that a described invention actually works. This paper surveys the general pattern common to these filings, argues that patent grant is being systematically misread as a form of scientific validation it was never designed to provide, and situates the recurring vagueness of these patents' operative mechanism sections against an entropic account of gravity in which no coherent local mechanism for the claimed effect could exist regardless of how the patent's language is arranged. The paper describes the legal standard of utility that governs what an examiner may and may not reject, the way that standard has been applied to perpetual motion machines, and the consequences for how a patent should be read as evidence. It is a survey of a pattern and not an audit of individual filings, and it does not assert that every patent in the area is unworkable by reason of its category alone.
\end{abstract}

\section{Introduction}

A recurring category of claim in antigravity discourse involves pointing to a granted patent, often filed decades earlier, as evidence that some agency with technical expertise has certified the underlying device as workable. Patents for gravity shielding devices, gravity wave generators, and inertia-reduction apparatus have indeed been granted in the United States, the United Kingdom, and elsewhere. Thomas Townsend Brown's British patent of 1928 and his American patent of 1965 are the classic instances, the latter describing an electrokinetic apparatus whose force is said to occur even in vacuum \cite{Brown1928,Brown1965}. The existence of these documents is frequently treated as more significant than it actually is.

\section{What Patent Examination Actually Certifies}

Patent examiners assess novelty, non-obviousness, and whether the application adequately describes the invention for a person skilled in the relevant art to construct it, called enablement. In the United States an invention must also have utility, and the guidance that governs examiners treats the question of utility as one of whether the claimed invention has a specific, substantial and credible use \cite{MPEP}. The standard is deliberately permissive. An examiner is expected to accept an asserted utility unless it is incredible in view of the established scientific understanding, and the burden of making that case rests on the office and not on the applicant. They do not, as a matter of practice or of law, verify that the described device functions as claimed.

The doctrine has teeth in a narrow band of cases. The best known concerns perpetual motion. In \emph{Newman v. Quigg}, the Federal Circuit affirmed the refusal of a patent on an energy machine claimed to produce more output than input, on the ground that an invention of that kind, contradicting established physical law, lacks operability and utility \cite{Newman1989}. The case is often read as showing that the office will reject claims that violate the laws of thermodynamics. It shows something narrower: that the office may do so when the violation is elementary and the applicant cannot produce a working device. Where the claim is framed in unfamiliar vocabulary, with field, resonance, polarization or coupling standing for a mechanism the examiner is not equipped to assess, the same patent passes. A granted patent, in short, certifies that an application was filed and processed correctly. It does not certify that the invention works.

\section{The Recurring Structure of These Filings}

Antigravity patents share a characteristic structure: a detailed, often genuinely careful description of the physical apparatus, its materials, its dimensions, and its assembly, paired with a comparatively vague account of why the described arrangement should produce the claimed gravitational effect, frequently invoking terms borrowed from legitimate physics, field, resonance, coupling, polarization, without specifying a mechanism precise enough to be tested or falsified. This asymmetry, rigorous engineering description paired with hand-waved physical justification, is exactly what one would expect from an inventor who has built something real, observed an ambiguous or null result, and filled the resulting explanatory gap with borrowed vocabulary rather than a specified causal chain, because the patent examination process does not require the gap to be filled with anything more rigorous than that.

The Brown patents illustrate the pattern with unusual clarity. The apparatus in each is described carefully, with electrodes, dielectrics and voltages, and the specification describes a force on the apparatus. Later work on asymmetric capacitors, discussed in the companion papers on the Biefeld-Brown effect and on lifters, found that the force in air is accounted for by ionic wind, while the claim that the force persists in vacuum has not survived controlled testing. The patent is a competent description of a device that produces a real effect, attached to an interpretation of the effect that the patent does not support. Reading the patent as an endorsement of the interpretation reverses its logical role.

\section{Why No Filling of the Gap Is Available}

On an account of gravitation as the large-scale bookkeeping of an ordered configuration's redistribution, tracked geometrically rather than transmitted as a locally sourced field, the vagueness common to these patents' mechanism sections is not a correctable drafting failure. There is no more precise mechanism available to specify, because the category of device being described, benchtop apparatus using magnets, capacitors, rotating masses, or resonant cavities, has no channel into the scale of ordered-configuration redistribution that produces measurable curvature in the first place. A more rigorously worded patent application describing the same devices would not close this gap; it would simply make the absence of a mechanism more visible, which is presumably why the vagueness persists.

The standard physics reaches a similar conclusion by a different route. Because gravitational and inertial mass are equal, to a precision now at a few parts in $10^{15}$, there is no body that is gravitationally coupled differently from another, and so no material can be a shield or a switch \cite{T2022}. A claim to switch gravity off for a body is a claim that the body can be made to violate this equality, and the equality has been tested with bodies of many compositions.

\section{How a Patent Should Be Read as Evidence}

A patent is evidence of a specific, limited set of facts: that a named person filed a description on a date, that the description was found to be new, and that the office did not find it incredible. A patent that describes a working device is evidence of the device only to the extent that the specification reports measurements, and the useful questions are the ones a reader would ask of any report. Does the specification give a measured force, with an uncertainty, under stated conditions? Was a control run in which the proposed active element was removed? Were alternative explanations excluded in the document itself? A specification that answers these is stronger evidence than a notice of grant. Most of the antigravity patents surveyed here answer none of them, which places them closer to a notebook entry than to a validation.

The patent literature also has a legitimate use in the study of such claims. Because patents are dated and public, they provide a timestamped record of what an inventor claimed to have built, and they are valuable to historians of the subject for exactly that reason. The criticism in this paper is aimed at the move from that record to an endorsement.

\section{Objections}

It might be argued that patent offices occasionally do get ahead of settled science, granting protection for mechanisms later vindicated, and that dismissing all such filings by category risks dismissing a genuine future exception along with the rest. This is true in the abstract but does not describe the antigravity patent record specifically, where the pattern across a century of filings has been consistent unverifiability rather than occasional vindication, and where the mechanism-vagueness described above is not a feature of early or unfashionable but eventually correct physics, it is a feature of applications describing devices with no available channel to the phenomenon they claim to affect.

A second objection is that a patent is a legal document and not a scientific one, so criticism of it on scientific grounds misses its purpose. The point is conceded, and it is the argument of this paper. The mistake lies with the reader who cites the document for a purpose it was never designed to serve.

\section{Conclusion}

A granted patent for a gravity-shielding or antigravity device demonstrates that a filing satisfied procedural requirements for novelty and description, not that the underlying physics works, and the persistent vagueness of these patents' causal mechanisms reflects a genuine absence of an available mechanism rather than an oversight correctable by better drafting. This paper surveys a pattern, tests no device and audits no individual claim, and its conclusion is limited to the evidentiary weight of a notice of grant.

\begin{thebibliography}{99}

\bibitem{Brown1928} T. T. Brown, A Method of Producing Force or Motion, British Patent 300,311, 1928.

\bibitem{Brown1965} T. T. Brown, Electrokinetic Apparatus, U.S. Patent 3,187,206, issued June 1, 1965.

\bibitem{MPEP} United States Patent and Trademark Office, \emph{Manual of Patent Examining Procedure}, section 2107, Guidelines for Examination of Applications for Compliance with the Utility Requirement.

\bibitem{Newman1989} \emph{Newman v. Quigg}, 877 F.2d 1575 (Fed. Cir. 1989).

\bibitem{T2022} P. Touboul et al., MICROSCOPE Mission: Final Results of the Test of the Equivalence Principle, \emph{Physical Review Letters} 129 (2022) 121102.

\end{thebibliography}

\end{document}
